Rest frequency $\nu_o = 1420.41$\,MHz

Detected frequency $\nu = 1421.23$\,MHz
\[
  \Delta\nu = \nu_o - \nu = 1420.41 - 1421.23 = -0.82\ \mathrm{MHz}
\]
\hfill(25\%)

Using Doppler's shift, the speed of the cloud is:
\[
  v = \frac{\Delta\nu}{\nu_o}\,c
    = \frac{-0.82}{1421.23}\left(3\times10^{10}\right)
    = -1.73\times10^{7}\ \mathrm{cm/s}
\]
\hfill(15\%)

Since $v$ is negative, then the cloud is moving toward us. \hfill(10\%)

(Since $\nu > \nu_o$, then the cloud is moving toward us)

If $\mathcal{M}$ is mass of black hole, $v$ is the speed of the cloud and $R$
is the orbital radius of cloud, then
\[
  \mathcal{M} = \frac{Rv^2}{G}
\]
\hfill(25\%)

$R = 0.2\,\mathrm{pc} = 0.2 \times (3.086 \times 10^{18}\,\mathrm{cm})
= 6.17 \times 10^{17}\,\mathrm{cm}$, and from the Table of Astronomical and
Physical constants: $G = 6.67 \times 10^{-8}\,\mathrm{cm}^3/\mathrm{s}^2\,\mathrm{g}$

Then,
\[
  \mathcal{M} = \frac{Rv^2}{G}
  = \frac{\left(6.17\times10^{17}\right)\left(-1.73\times10^{7}\right)^2}{6.67\times10^{-8}}
  = 2.78\times10^{39}\ \mathrm{gr} = 1.40\times10^{6}\ \mathcal{M}_\odot
\]
\hfill(25\%)
